\setcounter{section}{20}

  \zwischenueberschrift{Soal latihan}

\inputaufgabe
{}
{

Tentukan
\definitionsverweis {turunan eksterior}{}{}
dari
$1$-\definitionsverweis {bentuk diferensial}{}{}

\mavergleichskettedisp
{\vergleichskette
{ \omega
}
{ =} { \left(  x^2-y^3  \right) dx+x^3y^2dy
}
{ } {
}
{ } {
}
{ } {
}
}
{}{}{}
pada $\R^2$.

}
{} {}

\inputaufgabe
{}
{

Tentukan
\definitionsverweis {turunan eksterior}{}{}
dari
$1$-\definitionsverweis {bentuk diferensial}{}{}

\mavergleichskettedisp
{\vergleichskette
{ \omega
}
{ =} { xy^2dx+yzdy+x^3dz
}
{ } {
}
{ } {
}
{ } {
}
}
{}{}{}
pada $\R^3$.

}
{} {}

\inputaufgabegibtloesung
{}
{

Hitung
\definitionsverweis {turunan eksterior}{}{}
$d \omega$ dari
$1$-\definitionsverweis {bentuk diferensial}{}{}

\mavergleichskettedisp
{\vergleichskette
{ \omega
}
{ =} { { \frac{   x^2 }{  y } } dx- { \frac{   x  }{  y^2 } } dy
}
{ } {
}
{ } {
}
{ } {
}
}
{}{}{}
pada

\mavergleichskette
{\vergleichskette
{ U
}
{ = }{ \left\{ (x,y) \in \R^2  \mid   y \neq 0         \right\}
}
{  }{
}
{    }{
}
{   }{
}
}
{}{}{.}

}
{} {}

\inputaufgabegibtloesung
{}
{

Hitung turunan eksterior $d \omega$ dari bentuk diferensial

\mavergleichskettedisp
{\vergleichskette
{ \omega
}
{ =} {  e^{xz} dx \wedge dy - xyz dx \wedge dz + \left(  \sin ( \cos (xy) )  +y^{10}z^{100}  \right) dy \wedge dz
}
{ } {
}
{ } {
}
{ } {
}
}
{}{}{}
pada
\mathl{\R^3}{.}

}
{} {}

\inputaufgabegibtloesung
{}
{

Misalkan
\maabbeledisp {f} {\R} {\R
} { t } {f(t)
} {,}
diberikan oleh

\mavergleichskettedisp
{\vergleichskette
{ f(t)
}
{ =} { { \frac{   \sin^{ 3  } (t^4)  }{  1+t^2 } }
}
{ } {
}
{ } {
}
{ } {
}
}
{}{}{}

a) Hitung turunan eksterior dari $f$.

b) Hitung turunan eksterior dari $fdt$.

}
{} {}

\inputaufgabe
{}
{

Tentukan
\definitionsverweis {turunan eksterior}{}{} dari
$2$-\definitionsverweis {bentuk diferensial}{}{}

\mathdisp {\omega = xdx \wedge dy+xy^2zdy \wedge dz+xe^ydx\wedge dz} {  }
 pada $\R^3$.

}
{} {}

\inputaufgabegibtloesung
{}
{

Tunjukkan bahwa
$1$-\definitionsverweis {bentuk diferensial}{}{}

\mavergleichskette
{\vergleichskette
{ \omega
}
{ = }{ (2x- \sin  y )dx-x\cos  ydy
}
{  }{
}
{    }{
}
{   }{
}
}
{}{}{}
pada $\R^2$
\definitionsverweis {tertutup}{}{}
dan juga
\definitionsverweis {eksak}{}{.}

}
{} {}

\inputaufgabe
{}
{

Misalkan

\mavergleichskette
{\vergleichskette
{ U
}
{ \subseteq }{ \R^n
}
{  }{
}
{    }{
}
{   }{
}
}
{}{}{}
suatu
\definitionsverweis {himpunan bagian terbuka}{}{}
dan

\mavergleichskette
{\vergleichskette
{ \omega
}
{ = }{ \sum_{i = 1}^n f_i dx_i
}
{  }{
}
{    }{
}
{   }{
}
}
{}{}{}
suatu
$1$-\definitionsverweis {bentuk diferensial}{}{,}
yang dinyatakan dengan
\definitionsverweis {fungsi-fungsi yang terdiferensialkan secara kontinu}{}{}
\maabb {f_i} { U } { \R
} {}.
Tunjukkan bahwa $\omega$
\definitionsverweis {tertutup}{}{}
jika dan hanya jika

\mavergleichskettedisp
{\vergleichskette
{ \partial_i f_j
}
{ =} { \partial_j f_i
}
{ } {
}
{ } {
}
{ } {
}
}
{}{}{}
untuk semua $i,j$.

}
{} {}

\inputaufgabe
{}
{

Misalkan

\mavergleichskette
{\vergleichskette
{ U
}
{ \subseteq }{ \R^n
}
{  }{
}
{    }{
}
{   }{
}
}
{}{}{}
suatu
\definitionsverweis {himpunan bagian terbuka}{}{}
dan $\omega$ suatu
\definitionsverweis {bentuk diferensial}{}{} berderajat $1$ yang
\definitionsverweis {terdiferensialkan secara kontinu}{}{}
pada $U$, dengan
\definitionsverweis {medan vektor}{}{}
$F$ pada $U$ yang bersesuaian menurut
[[Riemannsche Mannigfaltigkeit/Vektorfelder und 1-Formen/Fakt|Lema 16.3]].
Tunjukkan bahwa $\omega$
\definitionsverweis {tertutup}{}{}
jika dan hanya jika $F$ memenuhi
\definitionsverweis {syarat keterintegralan}{}{.}

}
{} {}

\inputaufgabe
{}
{

Misalkan
\mathl{U \subseteq \R^n}{} suatu
\definitionsverweis {himpunan bagian terbuka}{}{}
dan $\omega$ suatu
\definitionsverweis {bentuk diferensial}{}{} berderajat $1$ yang
\definitionsverweis {terdiferensialkan secara kontinu}{}{}
pada $U$, dengan
\definitionsverweis {medan vektor}{}{}
$F$ pada $U$ yang bersesuaian menurut
[[Riemannsche Mannigfaltigkeit/Vektorfelder und 1-Formen/Fakt|Lema 16.3]].
Tunjukkan bahwa $\omega$
\definitionsverweis {eksak}{}{}
jika dan hanya jika $F$ merupakan
\definitionsverweis {medan gradien}{}{.}

}
{} {}

\inputaufgabe
{}
{

Misalkan $\omega$ suatu
\definitionsverweis {bentuk diferensial}{}{} \definitionsverweis {terdiferensialkan}{}{}
pada manifold $M$, dan
\maabb {\varphi} { L } { M
} {}
suatu pemetaan yang dua kali terdiferensialkan secara kontinu dan didefinisikan pada manifold $L$.

\aufzaehlungzwei {Misalkan $\omega$
\definitionsverweis {eksak}{}{.}
Tunjukkan bahwa bentuk diferensial tarik balik $\varphi^*\omega$ juga eksak.
} {Misalkan $\omega$
\definitionsverweis {tertutup}{}{.}
Tunjukkan bahwa
\definitionsverweis {bentuk diferensial tarik balik}{}{}
$\varphi^*\omega$ juga tertutup.
}

}
{} {}

Untuk soal-soal berikut, bandingkan
[[Vektorfeld/(x,y) nach (-y,x)/Nicht integrabel/Beispiel|Contoh 58.7 (Analisis (Osnabrück 2021–2023))]]
dan
[[Vektorfeld/Punktierte Ebene/(x,y) nach (-y,x) durch x^2+y^2/Integrabel nicht exakt/Beispiel|Contoh 58.10 (Analisis (Osnabrück 2021–2023))]].

\inputaufgabe
{}
{

Tunjukkan bahwa
$1$-\definitionsverweis {bentuk diferensial}{}{}

\mathdisp {-ydx+xdy} {  }

pada $\R^2$ tidak
\definitionsverweis {tertutup}{}{.}

}
{} {}

\inputaufgabe
{}
{

Tunjukkan bahwa
$1$-\definitionsverweis {bentuk diferensial}{}{}

\mathdisp {- { \frac{   y  }{  x^2+y^2 } }  dx+ { \frac{   x  }{  x^2+y^2 } } dy} {  }

pada $\R^2 \setminus \{0\}$
\definitionsverweis {tertutup}{}{,}
tetapi tidak
\definitionsverweis {eksak}{}{.}

}
{} {}

\inputaufgabegibtloesung
{}
{

Dalam dimensi positif, tunjukkan bahwa setiap
$n$-\definitionsverweis {bentuk diferensial}{}{}
kontinu $\omega$ pada suatu himpunan terbuka

\mavergleichskette
{\vergleichskette
{U
}
{ \subseteq }{\R^n
}
{  }{
}
{    }{
}
{   }{
}
}
{}{}{}
merupakan bentuk
\definitionsverweis {eksak}{}{.}

}
{} {}

\inputaufgabe
{}
{

Misalkan $M$ suatu
\definitionsverweis {manifold diferensiabel}{}{}
yang
\definitionsverweis {terhubung}{}{}
dan $\omega$ suatu
$1$-\definitionsverweis {bentuk diferensial}{}{}
\definitionsverweis {terdiferensialkan}{}{}
pada $M$. Tunjukkan bahwa $\omega$
\definitionsverweis {eksak}{}{}
jika dan hanya jika, untuk setiap lintasan
\definitionsverweis {terdiferensialkan secara kontinu}{}{}
\maabbdisp {\gamma} { [a,b] } { M
} {}
\definitionsverweis {integral lintasan}{}{}

\mathl{\int_\gamma \omega}{} hanya bergantung pada
\mathl{\gamma(a)}{} dan
\mathl{\gamma(b)}{}.

}
{} {}

\inputaufgabe
{}
{

Misalkan $M$ dan $N$ merupakan
\definitionsverweis {manifold diferensiabel}{}{}
dan
\maabbdisp {\varphi} { M } { N
} {}
suatu
\definitionsverweis {pemetaan yang dua kali terdiferensialkan secara kontinu}{}{.}
Misalkan $\omega$ suatu
$n$-\definitionsverweis {bentuk diferensial}{}{}
\definitionsverweis {terdiferensialkan}{}{}
pada $N$, dengan $n$ adalah dimensi $N$. Tunjukkan bahwa $\varphi^* \omega$ merupakan
\definitionsverweis {bentuk diferensial}{}{}
\definitionsverweis {tertutup}{}{}
pada $M$.

}
{} {}

\inputaufgabe
{}
{

Manakah di antara fungsi-fungsi berikut
\maabbdisp {} {\R_+} {\R
} {}
yang dapat diperluas menjadi fungsi
\definitionsverweis {diferensiabel}{}{}
pada
\definitionsverweis {titik batas}{}{}
$0$?
\aufzaehlungsechs{
\mathl{x^3+ \sin^{ 3  }  x -e^{-x}}{,}
}{
\mathl{{ \frac{   1  }{  x  } }}{,}
}{
\mathl{\sin  { \frac{   1  }{  x  } }}{,}
}{
\mathl{x \sin  { \frac{   1  }{  x  } }}{,}
}{
\mathl{{ \frac{   e^{ { \frac{   1  }{  x  } } }  }{  x  } }}{,}
}{
\mathl{x^2 \sin  { \frac{   1  }{  x  } }}{.}
}

}
{} {}

\inputaufgabe
{}
{

Misalkan
\mathl{H \subset \R^n}{} suatu
\definitionsverweis {setengah ruang}{}{.}
Misalkan
\mathl{Q\in H}{} suatu titik dan $Q \in U \subseteq H$, dengan $U$ suatu
\definitionsverweis {himpunan bagian terbuka}{}{}
dari $\R^n$. Tunjukkan bahwa $Q$ bukan titik batas $H$.

}
{} {}

\inputaufgabe
{}
{

Definisikan istilah \stichwort {difeomorfisme} {,} \stichwort {diferensial total} {} dan \stichwort {turunan tingkat tinggi} {} untuk
\definitionsverweis {setengah ruang}{}{}
\zusatzklammer {atau himpunan bagian terbuka darinya} {} {.}

}
{} {}

  \zwischenueberschrift{Soal untuk dikumpulkan}

\inputaufgabe
{3}
{

Tentukan
\definitionsverweis {turunan eksterior}{}{} dari
$1$-\definitionsverweis {bentuk diferensial}{}{}

\mathdisp {\omega = xy^2z^3dx+xyzdy+x^3yz^4dz} {  }
 pada $\R^3$.

}
{} {}

\inputaufgabe
{3}
{

Tentukan
\definitionsverweis {turunan eksterior}{}{} dari
$2$-\definitionsverweis {bentuk diferensial}{}{}

\mathdisp {\omega = xy^2dx \wedge dy+(x^3-y^2z^4)dy \wedge dz+ \sin (xy) dx \wedge dz} {  }

pada $\R^3$.

}
{} {}

\inputaufgabe
{5}
{

Misalkan

\mavergleichskette
{\vergleichskette
{ U
}
{ \subseteq }{ \R^n
}
{  }{
}
{    }{
}
{   }{
}
}
{}{}{}
\definitionsverweis {terbuka}{}{}
dan misalkan
\mathl{\omega_1 , \ldots , \omega_r}{} bentuk-bentuk diferensial pada $U$, dengan $\omega_i$ suatu
$k_i$-\definitionsverweis {bentuk diferensial}{}{.}
Temukan dan buktikan suatu rumus untuk

\mathdisp {d( \omega_1 \wedge \ldots \wedge \omega_r)} { . }

}
{} {}

\inputaufgabe
{5}
{

Tunjukkan bahwa
\definitionsverweis {bentuk diferensial}{}{}

\mavergleichskettedisp
{\vergleichskette
{ \omega
}
{ =} { \left(  2xy+3x^2-y e^{xy}  \right) dx + \left(  x^2-xe^{xy} +8y  \right)  dy
}
{ } {
}
{ } {
}
{ } {
}
}
{}{}{}
pada $\R^2$
\definitionsverweis {tertutup}{}{}
dan juga
\definitionsverweis {eksak}{}{.}

}
{} {}

\inputaufgabe
{6}
{

Tunjukkan bahwa
\definitionsverweis {setengah bidang}{}{}
$\R_{\geq 0} \times \R$ dan
\definitionsverweis {kuadran}{}{}

\mathl{\R_{\geq 0} \times \R_{\geq 0}}{}
\definitionsverweis {homeomorfik}{}{.}

}
{} {}

 [[Kategorie:Latexseite]]
